% ============================================================
% Chapter 4: Proposed Method: Specification and Design
%   Section: The Fact-Checking Pipeline
% Included from chapters/04-proposed-method/proposed-method.tex
% Source: docs/arquitectura-tecnica.md section 2.4 (seven stages),
% docs/decisions/retrieval.md, src/services/fact_checker/. The stage
% list of Section~\ref{subsec:discard-justification} follows this one.
% ============================================================

\section{The Fact-Checking Pipeline}
\label{sec:design-fact-checking}

Fact-checking is not one true-or-false step but a chain of seven stages
(Figure~\ref{fig:fact-checking}), each more expensive than the one
before it. Any stage can stop or downgrade the result, and no later
stage can undo a strong rejection by an earlier one. The governing
principle is to fail closed: when in doubt, the system leans towards
``not verified'' rather than towards a confident verdict that may be
wrong, because a false sense of certainty costs more than caution. The
only point where a language model takes part is bounded to a closed
classification that must cite its sources, and its output is then
checked by rules that do not depend on the model's own judgement.

\begin{figure}[ht]
\centering
\begin{adjustbox}{max width=\linewidth}
\begin{tikzpicture}[
    node distance=4.5mm,
    st/.style={draw, rounded corners=2pt, font=\footnotesize, align=left,
               text width=8.2cm, minimum height=7mm, fill=ganttplan!15,
               inner sep=3pt},
    outcome/.style={st, fill=black!6, text width=3.6cm, dashed},
    arr/.style={-{Stealth[length=2mm]}, thick},
    lbl/.style={font=\scriptsize\itshape, fill=white, inner sep=1pt,
                align=right, text width=2.2cm},
]
\node[st] (a) {\textbf{1. Admission} --- on topic, constructive enough, not a duplicate};
\node[st, below=of a] (b) {\textbf{2. Claim selection} --- two to four anchor claims};
\node[st, below=of b] (c) {\textbf{3. Evidence retrieval} --- three queries per claim, fused; own corpus};
\node[st, below=of c] (d) {\textbf{4. Ranking and pertinence gate} --- four factors; off-target sources cut};
\node[st, below=of d] (e) {\textbf{5. Verification} --- closed verdict, cited sources, verbatim quotes};
\node[st, below=of e] (f) {\textbf{6. Confidence recalibration} --- ungrounded verdicts downgraded};
\node[st, below=of f] (g) {\textbf{7. Aggregation} --- the article takes its worst claim's verdict};
\node[outcome, right=8mm of a] (x) {Rejected: no search, no model call; reasons stored};
\node[outcome, right=8mm of c] (y) {\texttt{searchUnavailable}: every query failed};
\node[outcome, right=8mm of e] (z) {\texttt{llmUnreachable}: the model never answered};
\draw[arr] (a) -- (b); \draw[arr] (b) -- (c); \draw[arr] (c) -- (d);
\draw[arr] (d) -- (e); \draw[arr] (e) -- (f); \draw[arr] (f) -- (g);
\draw[arr] (a) -- (x); \draw[arr] (c) -- (y); \draw[arr] (e) -- (z);
\draw[arr] (d.west) -- ++(-5mm,0) |- node[lbl, pos=0.25, left]{no evidence: model not asked} (f.west);
\end{tikzpicture}
\end{adjustbox}
\caption{The seven stages. Stages 2--6 run per claim, and claims run
concurrently; each claim records the furthest stage it reached. The
dashed boxes are outcomes reported apart from the verdicts.}
\label{fig:fact-checking}
\end{figure}

Stages~1 and~7 work on the article; stages~2 to~6 work on each claim,
in that order and never overlapped, since each feeds the next. A claim
whose verdict leaves stage~6 in doubt can, when the run enables it, go
back through stages~3 to~6 once more with a second, wider search
(Subsection~\ref{subsec:second-look}). What the
evaluation measures is the claim-checking part, stages~3 to~6, through
the same entry point the single-claim endpoint uses
(Section~\ref{sec:harness}).

% ============================================================
% Chapter 4: Proposed Method: Specification and Design
%   Section: The Fact-Checking Pipeline
%     Subsection: Admission Filter
% Included from chapters/04-proposed-method/06-fact-checking-pipeline/fact-checking-pipeline.tex
% Sources: docs/arquitectura-tecnica.md 2.4.1;
% backend/src/services/admission/README.md.
% ============================================================

\subsection{Admission Filter}
\label{subsec:admission}

Before anything is spent on verification, the article must pass three
independent checks. They judge the \emph{article} --- its topics, its
tone, its neighbours in the vector store --- and none needs evidence, a
search or a language model, which is why they come first and are the
cheapest stage.

\begin{itemize}
    \item \textbf{Topic relevance}: at least one topic's confidence
    must reach a minimum (0.35 by default). No topic has to dominate,
    since an article can legitimately touch several.
    \item \textbf{Positive impact}: a weighted average of
    constructiveness, inspirational value, hopefulness, positive
    sentiment, objectivity and social impact, in $[0,1]$, must reach a
    minimum (0.1 by default). Hard floors on objectivity and negative
    sentiment, which would reject an article whatever its average, exist
    behind a single switch and are off by default: the floors on
    constructiveness and inspiration rejected too many real articles,
    and now only lower the score and are recorded as reasons.
    \item \textbf{Not a duplicate}: the article's embedding is compared
    with its nearest stored neighbour. Below a relatedness threshold
    (0.80) the two are unrelated; between it and a duplicate threshold
    (0.90) they are \emph{related}, a useful signal that does not
    disqualify; above it the article is a duplicate. A stored article
    with the same URL never counts: it is the same article from an
    earlier run, and a re-analysis must not reject itself.
\end{itemize}

All three always run, even after one fails, so a rejection names every
reason. A rejected article is stored with its reasons and is never
searched or sent to a model. Every threshold can be overridden for one
run (Section~\ref{sec:dev-storage}).


% ============================================================
% Chapter 4: Proposed Method: Specification and Design
%   Section: The Fact-Checking Pipeline
%     Subsection: Claim Selection
% Included from chapters/04-proposed-method/06-fact-checking-pipeline/fact-checking-pipeline.tex
% Source: docs/arquitectura-tecnica.md 2.4.2; claim_selector.py.
% ============================================================

\subsection{Claim Selection}
\label{subsec:claim-selection}

Not every candidate claim of the enrichment stage is worth checking.
Two controls, both in the embedding space, reduce them to the few that
carry the article:

\begin{itemize}
    \item \textbf{Semantic deduplication}: when two candidates state the
    same fact in different words (similarity above 0.92), only one
    survives.
    \item \textbf{An anchor band}: the remaining candidates are ranked
    by how load-bearing they are for the article's thesis --- how
    central they are to its headline and lead, how much they share its
    main subjects, how specific they are --- and between two and four
    are kept. The ceiling stops an unusually dense article from
    multiplying the cost of verification; an article that yields fewer
    than the minimum is reported as such rather than padded.
\end{itemize}

Sentences that read as opinion or interpretation are dropped before the
ranking: no evidence can confirm or refute ``a hopeful step for the
sector''. The selector's precision --- the share of its proposals a
person accepts as checkable claims --- is measured by the daily
labelling batches, which ask the annotator to label or reject every
claim it proposes, with a reason (Section~\ref{sec:custom-dataset}).


% ============================================================
% Chapter 4: Proposed Method: Specification and Design
%   Section: The Fact-Checking Pipeline
%     Subsection: Evidence Retrieval
% Included from chapters/04-proposed-method/06-fact-checking-pipeline/fact-checking-pipeline.tex
% Source: docs/decisions/retrieval.md; retrieval/query_builder.py.
% ============================================================

\subsection{Evidence Retrieval}
\label{subsec:evidence-retrieval}

\paragraph{The failure the design prevents.} An early version built one
query per claim from its entities, figures and dates. A Spanish claim
that the Coyote's mail-order ACME purchases stood for post-war American
consumerism was searched as ``coyote ACME Estados Unidos Segunda
Guerra Mundial'', which every page explaining the Greek word
\emph{acme} matches; those pages were the best embedding match, the
model was shown three of them, and the claim came back \texttt{FALSE}
at 83\% confidence. Each stage had behaved as designed, and none could
produce the honest answer, which was \texttt{UNVERIFIED}.

\paragraph{Three queries, fused.} Each claim is now searched with a
planned set of three queries:

\begin{itemize}
    \item an \textbf{anchor} query --- entities, quoted figures and
    dates --- for precision on the subject;
    \item a \textbf{proposition} query --- the claim's own content
    words plus two anchors --- for what the claim asserts, which the
    anchor query lacks;
    \item a \textbf{refutation} query --- the anchors plus words that
    would contradict the claim, per language --- because every query
    built from a claim is phrased affirmatively and could otherwise only
    confirm it.
\end{itemize}

The three run concurrently and their ranked lists are fused by
reciprocal rank fusion \cite{cormack2009rrf}, which orders candidates by
how much the queries agree rather than by which query ran first, and
each candidate records which queries found it. When a claim lost its
subject on being separated from its article (``the programme trained
30 experts''), the headline words that the article's own keywords agree
on are restored to its queries, and the article's own URL is dropped
from the results, so an article cannot confirm itself.

\paragraph{Two sources of candidates.} The queries go to SearXNG, which
consults a closed list of engines measured one by one --- general web
engines and, at half weight, four academic indexes (arXiv, Crossref,
Semantic Scholar and PubMed) that keep answering when the web engines
rate-limit. In parallel, the system's own corpus of analysed articles is
searched by embedding, always excluding the article itself.

\paragraph{A funnel.} Reading a page in full is the most expensive
operation of the pipeline, so candidates are first pre-ranked cheaply
on their title and snippet, with the same semantic and lexical split
the ranking uses, and only the best eight are fetched through the
extraction cascade. Evidence pages never escalate to the headless
browser.

\paragraph{Search unavailable is not a verdict.} If every query of a
claim fails --- SearXNG unreachable, or answering empty while reporting
its engines down --- the claim is marked \texttt{searchUnavailable}
and stops here with a note that the web was never asked. Its verdict is
still \texttt{UNVERIFIED}, but it no longer looks like a judgement about
the claim, and the evaluation sets it apart. Before giving up, a query
SearXNG cannot answer is sent once to a second, keyless route.


% ============================================================
% Chapter 4: Proposed Method: Specification and Design
%   Section: The Fact-Checking Pipeline
%     Subsection: Evidence Ranking and the Pertinence Gate
% Included from chapters/04-proposed-method/06-fact-checking-pipeline/fact-checking-pipeline.tex
% Source: docs/decisions/retrieval.md; ranking/ranking_retrieval.py;
% settings.py (RANKING_*, EVIDENCE_MIN_PERTINENCE).
% ============================================================

\subsection{Evidence Ranking and the Pertinence Gate}
\label{subsec:ranking}

The fetched evidence is re-ranked by a score of four independent
factors:
\[
  \textit{relevance} = 0.45\cdot\textit{semantic}
    + 0.20\cdot\textit{lexical}
    + 0.20\cdot\textit{recency}
    + 0.15\cdot\textit{reliability}.
\]

\begin{itemize}
    \item \textbf{Semantic}: the embedding similarity between the claim
    and the evidence.
    \item \textbf{Lexical}: the share of the claim's own terms the
    evidence contains, anchors counted double. It is the counterweight
    to the semantic factor, which scores a page \emph{about the same
    subject} as highly as a page that addresses the claim --- exactly
    how the ACME etymology pages won.
    \item \textbf{Recency}: an exponential decay with a half-life of a
    year, so old evidence loses weight gradually rather than being
    discarded.
    \item \textbf{Reliability}: a rating per domain, from the source
    catalogue. Only the 37 configured domains have a real rating; any
    other domain gets a neutral 0.5 and is flagged as unrated, so the
    default is never presented as a rating.
\end{itemize}

The weights must sum to one, and the backend refuses to start if they
do not: the configuration template once shipped weights summing to 1.2,
and nothing noticed. The best five sources per claim are kept, at most
one per domain.

\paragraph{The pertinence gate.} In the same step, each source gets a
second score that leaves recency and reliability out:
\[
  \textit{pertinence} = 0.5\cdot\textit{semantic} + 0.5\cdot\textit{lexical}.
\]
A recent, reliable page about something else is still about something
else, and letting those two factors lift it would repeat the mistake
the gate exists to stop. A source whose pertinence is below a threshold
(0.25 by default) is cut \emph{before the model sees it}, with a reason
that names the figure: it cannot be cited, counted towards the evidence
minimum or treated as corroboration. The threshold is deliberately low
--- a floor against off-target retrieval, not a relevance ranking ---
because cutting real evidence turns a checkable claim into
\texttt{UNVERIFIED}, another way of being useless. It is a per-run
threshold, unlike the weights, because it is the knob that trades a
confident wrong answer for an honest one. Its default was reasoned, not
fitted: no labelled set existed when it was set.


% ============================================================
% Chapter 4: Proposed Method: Specification and Design
%   Section: The Fact-Checking Pipeline
%     Subsection: Verification with a Language Model
% Included from chapters/04-proposed-method/06-fact-checking-pipeline/fact-checking-pipeline.tex
% Source: docs/arquitectura-tecnica.md 2.4.5; verification/
% (per_source.py for the one-source-at-a-time rule).
% ============================================================

\subsection{Verification with a Language Model}
\label{subsec:llm-verification}

The model is given the claim, the article's headline and opening (so
that it reads the claim in its article's sense), and the ranked
evidence, numbered. It must return a structured answer:

\begin{itemize}
    \item a verdict from a closed set of five --- \texttt{TRUE},
    \texttt{PARTIALLY\_TRUE} (the central assertion holds, a detail does
    not), \texttt{MISLEADING} (technically true but deceptive, or
    sources of equal standing conflicting), \texttt{FALSE} and
    \texttt{UNVERIFIED};
    \item a confidence and a short explanation;
    \item the numbers of the evidence it used; and
    \item for each source, its stance --- \emph{supports},
    \emph{contradicts} or \emph{unrelated} --- with the verbatim passage
    it relied on.
\end{itemize}

It is told to answer \texttt{UNVERIFIED} when the evidence is empty or
does not bear on the claim, never to assert what the evidence does not
support, and to mark as \emph{unrelated} a source about a different
event even when it shares a place, a date or a figure with the claim.

Reducing the task to a closed classification with cited, quoted
evidence --- instead of a free-form verdict --- turns the question of
whether the model hallucinated into something that can be checked
mechanically in the next stage. Any quote that does not appear
verbatim in its source is discarded: an invented quotation is worse
than none. The model is served locally by default
(\texttt{llama3.2:3b}, a 3-billion-parameter model chosen for its
footprint, not by a benchmark), and any OpenAI-compatible provider can
take its place by configuration.

\paragraph{One source at a time.} Since 22~October the default asks
the model about each ranked source separately, in a short call: whether
that source supports the claim, contradicts it, is silent on its point
or is about something else, with the verbatim passage. The verdict then
follows a written rule rather than the model's judgement. Outlets are
counted, not pages --- five copies of one wire story are one source ---
and the side more outlets take wins; a tie is \texttt{UNVERIFIED}, and
so is a set of sources none of which takes a side. The confidence grows
with the share of outlets that agree, and the explanation is built from
the stances and their checked quotes, so that it cannot tell a different
story from them. This form never answers \texttt{PARTIALLY\_TRUE} or
\texttt{MISLEADING}. It was adopted on the results of
Section~\ref{sec:results-rq2}; the single call above remains a setting.

If the model cannot be reached after its retries, the claim is marked
\texttt{llmUnreachable} and reported apart from the verdicts, as a
failed search is. If there is no evidence at all, the model is not
asked: its answer would be discarded by the next stage anyway, and the
call is the slowest step of the pipeline.


% ============================================================
% Chapter 4: Proposed Method: Specification and Design
%   Section: The Fact-Checking Pipeline
%     Subsection: Confidence Recalibration
% Included from chapters/04-proposed-method/06-fact-checking-pipeline/fact-checking-pipeline.tex
% Source: docs/arquitectura-tecnica.md 2.4.6;
% verification/confidence_scorer.py.
% ============================================================

\subsection{Confidence Recalibration}
\label{subsec:recalibration}

This is the most important safeguard of the system, and the reason a
modest model can be used without compromising the reliability of the
result. Its rules do not depend on the model's own judgement of itself,
which is known to be poorly calibrated \cite{guo2017calibration}.

\begin{enumerate}
    \item \textbf{No evidence, no verdict.} Below the evidence minimum
    the verdict is forced to \texttt{UNVERIFIED} with confidence zero,
    and the model is never asked.
    \item \textbf{Confidence is recomputed.} The final confidence is
    $0.7$ times the model's confidence plus $0.3$ times an evidence
    quality score computed independently: the mean relevance of the
    retained evidence (half of it), the share of that evidence the model
    cited (three tenths) and the amount of evidence against the maximum
    expected (one fifth).
    \item \textbf{An ungrounded verdict is downgraded.} A definitive
    verdict that cites nothing, although evidence was available, becomes
    \texttt{UNVERIFIED} with its confidence capped at 0.4. Citing only
    sources the model itself called \emph{unrelated} counts as citing
    nothing.
    \item \textbf{A contradicted detail makes it partial.} A
    \texttt{TRUE} verdict while some source contradicts the claim
    becomes \texttt{PARTIALLY\_TRUE}: the central fact holds and a
    detail does not, which is the ordinary shape of a real verification.
    \item \textbf{One outlet is not confirmation.} A \texttt{TRUE} that
    rests on fewer than two independent domains keeps its verdict with
    its confidence capped at 0.6.
\end{enumerate}

The citation rules are independent of the model: they protect against
the hallucinations of any model, present or future, without retraining
or prompt changes. They were tested against a fake model that bluffs
in 83 ways --- citations out of range, repeated or of the wrong type,
fabricated or stitched quotes, verdicts outside the set, confidences of
5 or of infinity, malformed answers --- which found and fixed six ways
through (Section~\ref{sec:dev-testing}). Three remain, pinned as known
failures because each needs a policy rather than a fix: a source the
model calls both supporting and unrelated still grounds its citation; a
\texttt{FALSE} whose every cited source supports the claim stands; and
a \texttt{TRUE} whose every cited source contradicts it becomes only
\texttt{PARTIALLY\_TRUE}.


% ============================================================
% Chapter 4: Proposed Method: Specification and Design
%   Section: The Fact-Checking Pipeline
%     Subsection: The Second Look
% Included from chapters/04-proposed-method/06-fact-checking-pipeline/fact-checking-pipeline.tex
% Source: backend/src/services/fact_checker/second_look.py;
% docs/decisions/evaluation.md, "The second look".
% ============================================================

\subsection{The Second Look}
\label{subsec:second-look}

The first search for a claim is planned from the claim alone: who and
what it names, what it asserts, and its refutation
(Subsection~\ref{subsec:evidence-retrieval}). When those three queries
find nothing usable the claim ends \texttt{UNVERIFIED}; when they find
something thin, the model answers with low confidence. In both cases the
model's own explanation usually says what was missing --- ``no source
reports the 2024 figure'', ``the sources describe a different
programme''. The second look turns that sentence into a second search.

\begin{enumerate}
    \item \textbf{Follow-up queries.} The model is given the claim, the
    first verdict, its explanation and the titles of the sources it had,
    and asked for at most three queries, in the claim's language, aimed
    at what was missing: the original study or report, an official
    figure, the organisation's own announcement, an independent report
    of the same event. A model that writes none gets the claim itself as
    its query, a sentence the first search never sent whole.
    \item \textbf{Only new pages.} The queries go through the same
    search, with the same concurrency and fusion, and only pages the
    first pass never saw, on domains it did not already have, are kept;
    the article itself never is.
    \item \textbf{One judgement over both.} The first pass's evidence
    and the new pages are ranked, gated, judged by the model and
    recalibrated together, by the same code as the first pass. The
    second verdict replaces the first, which stays on the claim with the
    reason the second look was taken and the queries it sent, so a
    reader can see that the verdict came from a second search.
\end{enumerate}

It fires only when a run enables it, and only on a verdict in doubt:
\texttt{UNVERIFIED}, or a confidence below a threshold (0.5 by default).
It never fires on an outage --- with no search or no model, a second
search meets the same outage --- and never twice. It follows the
iterative retrieval of FIRE \cite{xie2025fire} and the question
generation of AVeriTeC's baseline \cite{schlichtkrull2023averitec},
reduced to one extra round: on a processor without a graphics card each
round costs a search, a handful of page fetches and two model calls.

Whether it should be on is not assumed. It is decided on the tuning
half of the custom set (Subsection~\ref{subsec:custom-split}) by the
tuning loop of Subsection~\ref{subsec:tuning-loop}, and measured on the
test half.


% ============================================================
% Chapter 4: Proposed Method: Specification and Design
%   Section: The Fact-Checking Pipeline
%     Subsection: Article-Level Aggregation
% Included from chapters/04-proposed-method/06-fact-checking-pipeline/fact-checking-pipeline.tex
% Source: docs/arquitectura-tecnica.md 2.4.7; fact_checker.py
% (_VERDICT_SEVERITY).
% ============================================================

\subsection{Article-Level Aggregation}
\label{subsec:aggregation}

An article's verdict is neither an average nor the verdict of its most
important claim: it is the worst of its claims' verdicts, in a fixed
order of severity,
\[
  \texttt{TRUE} < \texttt{PARTIALLY\_TRUE} < \texttt{UNVERIFIED} <
  \texttt{MISLEADING} < \texttt{FALSE}.
\]
\texttt{PARTIALLY\_TRUE} sits below \texttt{UNVERIFIED} because a
claim whose detail is off has at least been checked. The article's
confidence is the mean of its claims' confidences. An article is
publishable when it was admitted and its verdict is neither
\texttt{FALSE} nor \texttt{MISLEADING}: \texttt{UNVERIFIED} does not
block publication, because with a small local model it is the common
outcome, and the reader view shows each claim's own verdict.

\paragraph{Known limitation.} Because \texttt{UNVERIFIED} outranks
\texttt{TRUE}, a single claim without enough evidence decides the
article's verdict however many others were confirmed, and the mean
confidence mixes confidences of different verdicts. Separating ``how
much could be checked'' from ``what was found'' is left open
(Section~\ref{sec:limitations}). The reader view works around it for
the reader: a claim whose search or model failed reads ``not checked''
and is counted apart, and an article none of whose claims was judged is
shown as not checked rather than as unverified.

